\chapter{STL 常见函数}

\newenvironment{stlapi}{%
    \small
    \begin{longtable}{@{}%
        >{\raggedright\arraybackslash}p{0.35\textwidth}%
        >{\raggedright\arraybackslash}p{0.20\textwidth}%
        >{\raggedright\arraybackslash}p{0.38\textwidth}@{}}
    \toprule
    调用与参数 & 返回值 & 作用 \\
    \midrule
    \endhead
    \bottomrule
    \endlastfoot
}{\end{longtable}}

\section{常见容器}

容器对象记为 \texttt{a}，位置迭代器记为 \texttt{it}。共有接口 \texttt{size()} 返回元素数，\texttt{empty()} 返回是否为空。
支持迭代的容器用 \texttt{begin()} 和 \texttt{end()} 取得首元素与尾后迭代器。

\subsection{\texttt{vector} 与 \texttt{array}}

\begin{stlapi}
    \texttt{vector<T> a(n, x)} & 构造对象 & 建立含 \texttt{n} 个 \texttt{x} 的序列。 \\
    \texttt{array<T,N> a\{\}} & 构造对象 & 建立固定长度为 \texttt{N} 的数组，各元素值初始化。 \\
    \texttt{a[i]}、\texttt{a.at(i)} & 元素引用 & 访问下标 \texttt{i}。\texttt{at} 越界时抛出异常。 \\
    \texttt{a.front()}、\texttt{a.back()} & 元素引用 & 访问非空序列的首项或末项。 \\
    \texttt{a.push\_back(x)} & 无 & 在 \texttt{vector} 尾部添加 \texttt{x}。 \\
    \texttt{a.emplace\_back(args...)} & 新元素引用 & 在 \texttt{vector} 尾部用参数原地构造。 \\
    \texttt{a.pop\_back()} & 无 & 删除非空 \texttt{vector} 的末项。 \\
    \texttt{a.insert(it, x)} & 插入位置的迭代器 & 在 \texttt{vector} 的 \texttt{it} 之前插入 \texttt{x}。 \\
    \texttt{a.erase(it)} \newline
    \texttt{a.erase(f, l)} & 被删部分之后的迭代器 & 从 \texttt{vector} 删除单项或区间 \texttt{[f,l)}。 \\
    \texttt{a.resize(n, x)} & 无 & 改变 \texttt{vector} 长度，新增位置填 \texttt{x}。 \\
    \texttt{a.reserve(n)} & 无 & 为 \texttt{vector} 预留至少 \texttt{n} 项的容量，长度不变。 \\
    \texttt{a.assign(n, x)} & 无 & 将 \texttt{vector} 替换为 \texttt{n} 个 \texttt{x}。 \\
    \texttt{a.fill(x)} & 无 & 将 \texttt{array} 的全部元素置为 \texttt{x}。 \\
\end{stlapi}

\subsection{\texttt{string} 与 \texttt{string\_view}}

\begin{stlapi}
    \texttt{s.substr(pos, len)} & 新字符串或子视图 & 从 \texttt{pos} 取至多 \texttt{len} 个字符。省略长度时取到末尾。 \\
    \texttt{s.find(t, pos)} & 首个匹配下标或 \texttt{npos} & 从 \texttt{pos} 向后查找 \texttt{t}，起点默认是 $0$。 \\
    \texttt{s.rfind(t, pos)} & 最后匹配下标或 \texttt{npos} & 找起点不大于 \texttt{pos} 的最后一次匹配。 \\
    \texttt{s += t} & \texttt{string\&} & 在字符串尾部追加 \texttt{t}。 \\
    \texttt{s.erase(pos, len)} & \texttt{string\&} & 删除从 \texttt{pos} 开始的至多 \texttt{len} 个字符。 \\
    \texttt{s.replace(pos, len, t)} & \texttt{string\&} & 将指定子串替换成 \texttt{t}。 \\
    \texttt{stoi(s, \&pos, base)} \newline
    \texttt{stoll(s, \&pos, base)} & \texttt{int}、\texttt{ll} & 以 \texttt{base} 进制读取整数，\texttt{pos} 接收已读取字符数。默认十进制，位置参数可传空指针。 \\
    \texttt{to\_string(x)} & \texttt{string} & 将数值转成字符串。 \\
\end{stlapi}

\cppfile{code/stl/string_split.cpp}

\subsection{\texttt{stack}、\texttt{queue}、\texttt{deque} 与 \texttt{priority\_queue}}

\begin{stlapi}
    \texttt{a.push(x)} & 无 & 给栈、队列或优先队列插入 \texttt{x}。 \\
    \texttt{a.emplace(args...)} & 新元素引用或无 & 原地构造元素。采用默认底层容器时，栈和队列返回引用，优先队列无返回值。 \\
    \texttt{a.top()} & 元素引用 & 查看非空栈或优先队列的顶部。 \\
    \texttt{a.front()}、\texttt{a.back()} & 元素引用 & 查看非空队列或双端队列的首项或末项。 \\
    \texttt{a.pop()} & 无 & 删除栈顶、队首或堆顶。 \\
    \texttt{a.push\_front(x)} \newline
    \texttt{a.push\_back(x)} & 无 & 在双端队列指定端插入 \texttt{x}。 \\
    \texttt{a.pop\_front()} \newline
    \texttt{a.pop\_back()} & 无 & 从非空双端队列指定端删除一项。 \\
\end{stlapi}

\begin{cppcode}
priority_queue<ll> maxHeap;
priority_queue<ll, vector<ll>, greater<ll>> minHeap;
\end{cppcode}

\subsection{有序集合、映射与哈希容器}

\begin{stlapi}
    \texttt{a.insert(x)} & 迭代器或 \texttt{\{it,ok\}} & 插入键。映射容器传入 \texttt{\{key,value\}}。唯一键容器的 \texttt{ok} 表示是否新增。 \\
    \texttt{a.find(x)} & 迭代器 & 查找键 \texttt{x}，缺失时返回 \texttt{end()}。 \\
    \texttt{a.count(x)} & 匹配项数 & 统计等于 \texttt{x} 的键。 \\
    \texttt{a.lower\_bound(x)} \newline
    \texttt{a.upper\_bound(x)} & 迭代器 & 在有序容器内分别找首个 $\ge x$、首个 $>x$ 的键，默认按升序。 \\
    \texttt{a.equal\_range(x)} & 一对迭代器 & 返回与键 \texttt{x} 等价的半开区间。 \\
    \texttt{a.erase(x)}、\texttt{a.erase(it)} & 删除数量、下一迭代器 & 按键删除全部匹配项，或按迭代器删除一项。 \\
    \texttt{a[key]} & 映射值引用 & 访问 \texttt{map} 或 \texttt{unordered\_map}。缺键时插入值初始化的元素。 \\
    \texttt{a.at(key)} & 映射值引用 & 查找已有键，缺键时抛出异常。 \\
    \texttt{a.reserve(n)} & 无 & 按负载因子为哈希容器预留可容纳 \texttt{n} 项的桶。 \\
    \texttt{a.clear()} & 无 & 删除关联容器的全部元素。 \\
\end{stlapi}

\subsection{\texttt{bitset}、\texttt{pair} 与 \texttt{tuple}}

\begin{stlapi}
    \texttt{b.test(i)} & \texttt{bool} & 检查第 \texttt{i} 位，从最低位编号，越界时抛出异常。 \\
    \texttt{b.set(i, v)}、\texttt{b.reset(i)} \newline
    \texttt{b.flip(i)} & 位集引用 & 分别赋值、清零、取反。省略下标时可对全体位操作。 \\
    \texttt{b.count()} & 置位数 & 统计值为 $1$ 的位。 \\
    \texttt{b.any()}、\texttt{b.none()} \newline
    \texttt{b.all()} & \texttt{bool} & 分别判断是否存在置位、是否全零、是否全为一。 \\
    \texttt{p.first}、\texttt{p.second} & 成员 & 访问 \texttt{pair} 的两项。 \\
    \texttt{get<I>(t)} & 第 \texttt{I} 项的引用 & 按从 $0$ 开始的编号访问 \texttt{pair} 或 \texttt{tuple}。 \\
    \texttt{make\_pair(x,y)} \newline
    \texttt{make\_tuple(args...)} & 二元组、元组 & 从传入值推导类型并构造对象。 \\
\end{stlapi}

\section{排序、二分与查找}

函数区间用迭代器 \texttt{[f,l)} 表示。带 \texttt{cmp} 的排序和二分调用使用同一比较关系。
二分查找与集合运算的输入区间须按该关系有序。

\begin{stlapi}
    \texttt{sort(f,l,cmp)} \newline
    \texttt{stable\_sort(f,l,cmp)} & 无 & 按 \texttt{cmp} 排序。后者保留等价项的原顺序。省略比较器时按默认顺序。 \\
    \texttt{nth\_element(f,nth,l)} & 无 & 把排序后应位于 \texttt{nth} 的元素放到该处，两侧分别满足划分关系。 \\
    \texttt{partial\_sort(f,m,l)} & 无 & 将全部元素中最小的 \texttt{m-f} 项排序后放在前段。 \\
    \texttt{lower\_bound(f,l,x)} \newline
    \texttt{upper\_bound(f,l,x)} & 迭代器 & 默认升序下分别找首个 $\ge x$、首个 $>x$ 的位置，找不到返回 \texttt{l}。 \\
    \texttt{equal\_range(f,l,x)} & 一对迭代器 & 同时取得上下界，返回和 \texttt{x} 等价的区间。 \\
    \texttt{binary\_search(f,l,x)} & \texttt{bool} & 查询有序区间中是否有等价元素。 \\
    \texttt{find(f,l,x)} \newline
    \texttt{find\_if(f,l,pred)} & 迭代器 & 找首个等于 \texttt{x} 或满足 \texttt{pred} 的元素，找不到返回 \texttt{l}。 \\
    \texttt{count(f,l,x)} \newline
    \texttt{count\_if(f,l,pred)} & 匹配数量 & 统计等于 \texttt{x} 或满足 \texttt{pred} 的元素。 \\
    \texttt{unique(f,l)} & 新的逻辑尾 & 合并相邻等价元素。对容器调用 \texttt{erase} 移除逻辑尾之后的元素。 \\
    \texttt{remove(f,l,x)} \newline
    \texttt{remove\_if(f,l,pred)} & 新的逻辑尾 & 保留不等于 \texttt{x} 或不满足 \texttt{pred} 的元素及其相对顺序，之后可调用 \texttt{erase}。 \\
\end{stlapi}

\cppfile{code/stl/sorted_search.cpp}
\cppfile{code/stl/kth_smallest.cpp}

\begin{cppcode}
sort(a.begin(), a.end());
a.erase(unique(a.begin(), a.end()), a.end());
\end{cppcode}

\section{数值与序列函数}

\begin{stlapi}
    \texttt{gcd(a,b)}、\texttt{lcm(a,b)} & 两实参的公共整数类型 & 求最大公约数、最小公倍数。参数绝对值和结果须可表示。 \\
    \texttt{accumulate(f,l,init)} & 初值的类型 & 从 \texttt{init} 开始累加。传 \texttt{0LL} 可按 \texttt{ll} 累加。 \\
    \texttt{iota(f,l,x)} & 无 & 从 \texttt{x} 开始逐项赋值，每次将当前值加一。 \\
    \texttt{partial\_sum(f,l,out)} & 输出尾后迭代器 & 将各前缀的累加结果写入 \texttt{out}，累加类型取自输入元素。 \\
    \texttt{min\_element(f,l)} \newline
    \texttt{max\_element(f,l)} & 迭代器 & 返回首个最小或最大元素，空区间返回 \texttt{l}。 \\
    \texttt{min(a,b)}、\texttt{max(a,b)} & 选中参数的常量引用 & 返回两值中的较小值或较大值。 \\
    \texttt{clamp(x,lo,hi)} & 选中参数的常量引用 & 将值限制在闭区间 \texttt{[lo,hi]}，要求 \texttt{lo <= hi}。 \\
    \texttt{fill(f,l,x)}、\texttt{swap(a,b)} & 无 & 填充区间或交换两个对象。 \\
    \texttt{copy(f,l,out)} \newline
    \texttt{transform(f,l,out,op)} & 输出尾后迭代器 & 将原值或逐项调用 \texttt{op} 的结果写到从 \texttt{out} 开始的位置。 \\
    \texttt{reverse(f,l)} & 无 & 反转指定区间。 \\
    \texttt{rotate(f,m,l)} & 旋转后的分界迭代器 & 交换前段 \texttt{[f,m)} 与后段 \texttt{[m,l)} 的位置。 \\
    \texttt{next\_permutation(f,l)} \newline
    \texttt{prev\_permutation(f,l)} & \texttt{bool} & 转到字典序下一项或上一项。到达端点后返回假，并回到另一端。 \\
    \texttt{merge(f,l,g,h,out)} & 输出尾后迭代器 & 将有序区间 \texttt{[f,l)}、\texttt{[g,h)} 归并到从 \texttt{out} 开始的区间。 \\
    \texttt{set\_intersection} \newline
    \texttt{set\_union}、\texttt{set\_difference} & 输出尾后迭代器 & 调用参数同 \texttt{merge}，分别输出交集、并集与第一段减第二段的差集。 \\
\end{stlapi}

\cppfile{code/stl/permutations.cpp}

\begin{cppcode}
vector<int> common;
set_intersection(a.begin(), a.end(), b.begin(), b.end(),
                 back_inserter(common));
\end{cppcode}

\section{竞赛常用工具}

\subsection{比较器、迭代器与随机数}

\begin{stlapi}
    \texttt{[captures](args)\{...\}} & Lambda 对象 & \texttt{sort} 的比较器返回左项是否严格在右项之前。 \\
    \texttt{auto\& [x, y] = p} & 引用绑定 & 分别绑定到 \texttt{pair} 或二元组的两项。去掉 \texttt{\&} 时按值绑定。 \\
    \texttt{next(it, k)}、\texttt{prev(it, k)} & 移动后的迭代器 & 从 \texttt{it} 前进或后退 \texttt{k} 步，默认一步，原迭代器不变。 \\
    \texttt{distance(f,l)} & 迭代器差值类型 & 计算从 \texttt{f} 到 \texttt{l} 的距离。 \\
    \texttt{back\_inserter(a)} & 输出迭代器 & 每写入一项便向容器尾部追加一项。 \\
    \texttt{mt19937 rng(seed)} & 生成器对象 & 用 \texttt{seed} 构造生成器，通过 \texttt{rng()} 取一个值。 \\
    \texttt{uniform\_int\_distribution<T>} \newline
    \texttt{dist(lo,hi)} & 分布对象 & 构造整数分布，调用 \texttt{dist(rng)} 从闭区间取值。 \\
    \texttt{shuffle(f,l,rng)} & 无 & 用 \texttt{rng} 打乱随机访问区间。 \\
\end{stlapi}

\subsection{整数哈希、位运算与 128 位整数}

\begin{stlapi}
    \texttt{\_\_builtin\_popcountll(x)} & \texttt{int} & 返回 64 位掩码的置位数。 \\
    \texttt{\_\_builtin\_ctzll(x)} \newline
    \texttt{\_\_builtin\_clzll(x)} & \texttt{int} & 返回尾零或前导零的个数，参数须非零。 \\
    \texttt{toString(x)} & \texttt{string} & 将 \texttt{lll} 转成十进制字符串。 \\
    \texttt{parseInt128(s)} & \texttt{optional<lll>} & 将字符串解析为 \texttt{lll}，失败时返回空值。 \\
\end{stlapi}

\begin{cppcode}
unordered_map<ll, int, SplitMix64Hash> freq;
\end{cppcode}

\cppfile{code/stl/splitmix_hash.hpp}
\cppfile{code/stl/bit_ops.cpp}
\cppfile{code/stl/int128_io.hpp}

\subsection{PBDS 有序树（GCC 扩展）}

\begin{stlapi}
    \texttt{a.order\_of\_key(x)} & 元素数量 & 统计严格小于 \texttt{x} 的键。 \\
    \texttt{a.find\_by\_order(k)} & 迭代器 & 访问第 \texttt{k} 小的键，从 $0$ 开始，越界返回 \texttt{end()}。 \\
    \texttt{ms.insert(x)}、\texttt{ms.erase(x)} & 无、\texttt{bool} & 以下封装允许重复值。插入一项或删除一项，删除成功时返回真。 \\
    \texttt{ms.countLess(x)} & 元素数量 & 用下列封装统计严格小于 \texttt{x} 的项，计入重复值。 \\
    \texttt{ms.kth(k)} & \texttt{optional<ll>} & 返回从 $0$ 开始的第 \texttt{k} 小值，越界返回空值。 \\
\end{stlapi}

\cppfile[firstline=4]{code/stl/pbds_ordered_tree.hpp}
